\subsection{Cohomology of a rational-curve fibration}
\label{subsec:p1-cohomology}

The two remaining cases use integration along a smooth
$\mathbb P^1$-fibration. The following lemma expresses this operation
through an involution and a product rule, without requiring the
fibration to be the projectivization of a vector bundle.

\begin{lemma}\label{lem:p1-involution}
Let $p:B\to D$ be a smooth projective morphism of smooth connected
projective varieties whose fibres are isomorphic to $\mathbb P^1$.
Set $c=c_1(T_p)$ and $\delta=p^*p_*$.
Every class on $B$ is uniquely $a+cb$, with $a,b$ pulled back from $D$,
and $c^2$ is pulled back from $D$. Consequently
\[
 \overline{a+cb}=a-cb
\]
defines a ring involution. On even cohomology,
\begin{align}
 \delta(a+cb)&=2b,&
 \bar\alpha&=\alpha-c\delta\alpha,\label{eq:p1-involution}\\
 \delta(\alpha\gamma)
   &=(\delta\alpha)\gamma+\bar\alpha\,\delta\gamma.
   \label{eq:p1-product}
\end{align}
For a class $h$ of weight one with fibre degree one, both $2h-c$ and
$\alpha-h\delta\alpha$ are pulled back from $D$.
If $\alpha$ has top degree on $B$, then
$\int_B\bar\alpha=-\int_B\alpha$.
\end{lemma}

\begin{proof}
The restrictions of $1$ and $c/2$ form a basis of the cohomology of every
fibre. Leray--Hirsch therefore gives the asserted unique decomposition,
and fibre integration gives $p_*c=2$ and $p_*p^*a=0$.
The higher direct images of $\mathcal O_B$ are $p_*\mathcal O_B=\mathcal
O_D$ and $R^jp_*\mathcal O_B=0$ for $j>0$. Relative
Grothendieck--Riemann--Roch~\cite{BorelSerre1958} consequently gives
\[
 p_*\operatorname{td}(T_p)=1.
\]
Its component of weight one on $D$ is $p_*c^2/12=0$.
Writing $c^2=a+cb$ now gives $2b=0$, so $c^2$ is pulled back. This proves
that negating $c$ preserves products. The first two identities follow
from the decomposition, and multiplication of two such decompositions
proves~\eqref{eq:p1-product}.

Write $h=a_0+c/2$. Then $2h-c=2a_0$, while
$\alpha-h\delta\alpha=a-2a_0b$ for $\alpha=a+cb$.
Finally, in top degree the summand pulled back from $D$ vanishes by
dimension. The involution negates the remaining summand.
\end{proof}

In particular, if $\bar h=\ell-h$ and $\ell=2h-c$, repeated use of
\eqref{eq:p1-product} gives the polynomial identities
\begin{equation}\label{eq:p1-powers}
 \delta(h^a)=\sum_{k=0}^{a-1}h^k(\ell-h)^{a-1-k},\qquad
 \delta(c^a)=(1-(-1)^a)c^{a-1}\quad(a\ge1).
\end{equation}
These formulas are obtained by induction, so do not require division
by a cohomology class.
